\section{Solved Problems in Simple Linear Regression}

\subsection*{Problem Statements}

% Remember
\begin{problema}
\label{prob:ab5bf4a}
The relationship between study hours ($X$) and the grade obtained ($Y$) on a modular statistics exam is studied. Random data were collected from $n=10$ students:
\begin{center}
\begin{tabular}{|l|c|c|c|c|c|c|c|c|c|c|}\hline
\textbf{Hours ($X$)} & 2 & 3 & 4 & 5 & 6 & 7 & 8 & 9 & 10 & 12 \\ \hline
\textbf{Grade ($Y$)} & 45 & 50 & 55 & 60 & 65 & 70 & 75 & 80 & 85 & 90 \\ \hline
\end{tabular}
\end{center}
\begin{enumerate}
\item Calculate the ordinary least-squares (OLS) estimators of the intercept ($\hat{\beta}_0$) and slope ($\hat{\beta}_1$).
\item Write the empirical equation of the fitted linear regression model $\hat Y=\hat\beta_0+\hat\beta_1X$.
\item Predict the expected grade for a student who studies 7.5 hours before the exam.
\item Calculate $R^2$ and interpret it in the educational context.
\end{enumerate}
\end{problema}

% Understand
\begin{problema}
\label{prob:e5cf043}
The structural relationship between postal-package weight ($X$, in kg) and final shipping cost ($Y$, in dollars) is analyzed for a logistics company. Statistical software reports, from $n=20$ shipments with mean mass $\bar X=10.0$ kg,
\begin{align}
\hat Y=5.0+2.5X,\quad R^2=0.85,\quad \text{SST}=200.0.
\end{align}
\begin{enumerate}
\item Interpret the numerical estimates of the intercept ($\hat\beta_0=5.0$) and slope ($\hat\beta_1=2.5$).
\item Calculate the point prediction for a package weighing exactly 15.0 kg.
\item From $\text{SST}=200.0$ and $R^2$, determine the Regression Sum of Squares ($\text{SSR}$) and Residual Error Sum of Squares ($\text{SSE}$).
\item Calculate the sample variance of the estimation error $\hat\sigma^2$ and the model residual standard error ($S_e$).
\end{enumerate}
\end{problema}

% Apply
\begin{problema}
\label{prob:5a642fc}
A power-grid engineer analyzes whether fluctuations in outdoor ambient temperature ($X$, in $^\circ$C) have a statistically demonstrable impact on total daily energy consumption ($Y$, in megawatt-hours, MWh). For $n=15$ summer days, the descriptive metrics are
\begin{align}
\bar X=25.0\text{ }^\circ\text{C},\quad \bar Y=450.0\text{ MWh},\quad S_{XX}=500.0,\quad S_{XY}=2000.0.
\end{align}
\begin{enumerate}
\item Calculate the OLS estimates $\hat\beta_0$ and $\hat\beta_1$ and write the regression equation.
\item If the total sum of squares for electricity consumption is $\text{SST}=10000.0$, calculate the Regression Sum of Squares ($\text{SSR}$) and Residual Error Sum of Squares ($\text{SSE}$).
\item Test slope significance ($H_0:\beta_1=0$ versus $H_a:\beta_1\ne0$) at $\alpha=0.05$, calculating $t=\hat\beta_1/SE(\hat\beta_1)$ where $SE(\hat\beta_1)=\sqrt{\text{MSE}/S_{XX}}$, and compare with $t_{0.025,13}=2.160$.
\end{enumerate}
\end{problema}

% Analyze
\begin{problema}
\label{prob:e97453b}
\textbf{Analytical derivation of Gauss's normal equations and exact OLS unbiasedness:} consider the population linear regression model $Y_i=\beta_0+\beta_1X_i+\varepsilon_i$, $i=1,\ldots,n$, where $X_i$ is nonrandom with $S_{XX}>0$, and errors satisfy $E[\varepsilon_i]=0$, $\text{Var}(\varepsilon_i)=\sigma^2$, and $\text{Cov}(\varepsilon_i,\varepsilon_j)=0$ for $i\ne j$.
\begin{enumerate}
\item Defining $Q(\beta_0,\beta_1)=\sum_{i=1}^n(Y_i-\beta_0-\beta_1X_i)^2$, apply the first-order conditions ($\partial Q/\partial\beta_0=0$, $\partial Q/\partial\beta_1=0$) to derive Gauss's \textbf{normal equations}.
\item Solve the normal system to obtain $\hat\beta_1=\dfrac{\sum(X_i-\bar X)(Y_i-\bar Y)}{\sum(X_i-\bar X)^2}$ and $\hat\beta_0=\bar Y-\hat\beta_1\bar X$.
\item Using linearity of expectation and substituting the true model for $Y_i$, prove that $\hat\beta_1$ and $\hat\beta_0$ are strictly \textbf{unbiased} estimators.
\end{enumerate}
\end{problema}

% Evaluate
\begin{problema}
\label{prob:5c55133}
A metallurgical engineer studies cyclic-fatigue strength in megacycles ($Y$) of a structural alloy as a function of added titanium percentage ($X$). The model fitted to $n=18$ specimens is $\hat Y=120.0+15.0X$, with $\bar X=3.0\%$, $S_{XX}=40.0$, and $S_e=4.2$ megacycles.
\begin{enumerate}
\item Calculate the standard error of the population mean response $\text{SE}(\hat\mu_{Y|X_0})=S_e\sqrt{1/n+(X_0-\bar X)^2/S_{XX}}$ and the individual prediction standard error $\text{SE}(Y_{\text{pred}|X_0})=S_e\sqrt{1+1/n+(X_0-\bar X)^2/S_{XX}}$, both at $X_0=6.0\%$ ($t_{0.025,16}=2.120$).
\item Construct the 95\% confidence interval for the mean response and the 95\% prediction interval for an individual specimen at $X_0=6.0\%$. Explain why the prediction interval is substantially wider, even as $n\to\infty$.
\end{enumerate}
\end{problema}

% Create
\begin{problema}
\label{prob:58d3723}
Design an original simple-linear-regression problem (different from grades, packages, energy, alloys, automobiles, or advertising): a delivery application models delivery time ($Y$, in minutes) as a function of distance traveled ($X$, in km). For $n=12$ orders, summarize $\bar X=5.0$ km, $\bar Y=22.0$ min, $S_{XY}=54.0$, $S_{XX}=30.0$, and $\text{SST}=150.0$. Fit the OLS model, predict delivery time for an 8-km order, and calculate and interpret $R^2$.
\end{problema}

\subsection*{Detailed Solutions}

% Remember
\begin{solproblema}[prob:ab5bf4a]
\begin{enumerate}
\item $\bar X=66/10=6.6$, $\bar Y=675/10=67.5$. With $\sum X_iY_i=4895.0$ and $\sum X_i^2=508.0$:
\begin{align}
S_{XY}=4895.0-10(6.6)(67.5)=440.0,\qquad S_{XX}=508.0-10(6.6)^2=72.4.
\end{align}
\begin{align}
\hat\beta_1=\frac{440.0}{72.4}\approx6.0773,\qquad \hat\beta_0=67.5-(6.0773)(6.6)\approx27.3898.
\end{align}
\item $\hat Y=27.39+6.08X$.
\item For $X=7.5$: $\hat Y=27.3898+6.0773(7.5)\approx72.97$ points.
\item $\text{SSR}=\hat\beta_1S_{XY}=(6.0773)(440.0)\approx2674.03$. With $\sum Y_i^2=46825.0$: $\text{SST}=46825.0-10(67.5)^2=1262.5$. Thus $R^2=2674.03/1262.5\approx0.9972$ (99.72\%): 99.72\% of grade variability is explained by study hours, an almost functional linear relationship.
\end{enumerate}
\end{solproblema}

% Understand
\begin{solproblema}[prob:e5cf043]
\begin{enumerate}
\item $\hat\beta_0=5.0$: the fixed base processing cost, independent of weight ($X=0$), is \$5.00. $\hat\beta_1=2.5$: each additional kilogram increases shipping cost by \$2.50.
\item For $X=15.0$: $\hat Y=5.0+2.5(15.0)=42.50$ dollars.
\item $\text{SSR}=R^2\times\text{SST}=(0.85)(200.0)=170.0$. $\text{SSE}=\text{SST}-\text{SSR}=200.0-170.0=30.0$.
\item With $\nu_E=20-2=18$: $\hat\sigma^2=30.0/18\approx1.6667$, $S_e=\sqrt{1.6667}\approx1.2910$ dollars.
\end{enumerate}
\end{solproblema}

% Apply
\begin{solproblema}[prob:5a642fc]
\begin{enumerate}
\item $\hat\beta_1=2000.0/500.0=4.0$ MWh/$^\circ$C. $\hat\beta_0=450.0-4.0(25.0)=350.0$. Thus $\hat Y=350.0+4.0X$.
\item $\text{SSR}=\hat\beta_1S_{XY}=(4.0)(2000.0)=8000.0$. $\text{SSE}=10000.0-8000.0=2000.0$. Hence $R^2=0.800$.
\item With $\nu_E=13$, $\text{MSE}=2000.0/13\approx153.85$, $SE(\hat\beta_1)=\sqrt{153.85/500.0}\approx0.5547$, and $t=4.0/0.5547\approx7.211$. Since $|t|=7.211\gg2.160$, \textbf{reject $H_0$}: temperature and electricity consumption have a strongly significant dependence.
\end{enumerate}
\end{solproblema}

% Analyze
\begin{solproblema}[prob:e97453b]
\begin{enumerate}
\item Differentiating $Q(\beta_0,\beta_1)=\sum(Y_i-\beta_0-\beta_1X_i)^2$ and setting derivatives to zero:
\begin{align}
\frac{\partial Q}{\partial\beta_0}=-2\sum(Y_i-\beta_0-\beta_1X_i)=0,\qquad \frac{\partial Q}{\partial\beta_1}=-2\sum X_i(Y_i-\beta_0-\beta_1X_i)=0,
\end{align}
which yield Gauss's \textbf{normal equations}: $n\hat\beta_0+\hat\beta_1\sum X_i=\sum Y_i$ and $\hat\beta_0\sum X_i+\hat\beta_1\sum X_i^2=\sum X_iY_i$.
\item From the first equation, $\hat\beta_0=\bar Y-\hat\beta_1\bar X$. Substitution into the second and solving gives
\begin{align}
\hat\beta_1\left[\sum X_i^2-n\bar X^2\right]=\sum X_iY_i-n\bar X\bar Y\implies\hat\beta_1=\frac{\sum(X_i-\bar X)(Y_i-\bar Y)}{\sum(X_i-\bar X)^2}.
\end{align}
\item Substituting $Y_i=\beta_0+\beta_1X_i+\varepsilon_i$ and using $\sum(X_i-\bar X)=0$ and $\sum(X_i-\bar X)X_i=S_{XX}$:
\begin{align}
\hat\beta_1=\beta_1+\sum_{i=1}^n\left(\frac{X_i-\bar X}{S_{XX}}\right)\varepsilon_i.
\end{align}
Taking expectations with $E[\varepsilon_i]=0$ gives $E[\hat\beta_1]=\beta_1$. For the intercept, $E[\hat\beta_0]=E[\bar Y]-\bar X E[\hat\beta_1]=(\beta_0+\beta_1\bar X)-\bar X\beta_1=\beta_0$. Both estimators are unbiased.
\end{enumerate}
\end{solproblema}

% Evaluate
\begin{solproblema}[prob:5c55133]
\begin{enumerate}
\item At $X_0=6.0\%$: $\text{SE}(\hat\mu_{Y|X_0})=4.2\sqrt{1/18+(3.0)^2/40.0}=4.2\sqrt{0.28056}\approx2.2246$ megacycles. $\text{SE}(Y_{\text{pred}|X_0})=4.2\sqrt{1.28056}\approx4.7528$ megacycles.
\item With $\hat Y_0=120.0+15.0(6.0)=210.0$ and $t_{0.025,16}=2.120$,
\begin{align}
CI_{95\%}(\mu_{Y|X_0=6})&=210.0\pm2.120(2.2246)=[205.28,\;214.72],\\
PI_{95\%}(Y_{\text{pred}|X_0=6})&=210.0\pm2.120(4.7528)=[199.92,\;220.08].
\end{align}
The prediction interval ($\pm10.08$) is more than twice as wide as the mean-response confidence interval ($\pm4.72$). The confidence interval estimates the mean of infinitely many specimens under the same conditions and collapses toward zero width as $n\to\infty$, whereas a prediction interval must include the variance of one future observation, so it cannot narrow below $\pm t_{\alpha/2}\sigma$.
\end{enumerate}
\end{solproblema}

% Create
\begin{solproblema}[prob:58d3723]
\begin{align}
\hat\beta_1=\frac{S_{XY}}{S_{XX}}=\frac{54.0}{30.0}=1.8\text{ min/km},\qquad \hat\beta_0=22.0-1.8(5.0)=13.0\text{ min}.
\end{align}
The fitted model is $\hat Y=13.0+1.8X$. For an $X=8$ km order, $\hat Y=13.0+1.8(8)=27.4$ minutes. The regression sum of squares is $\text{SSR}=\hat\beta_1S_{XY}=(1.8)(54.0)=97.2$, so
\begin{align}
R^2=\frac{\text{SSR}}{\text{SST}}=\frac{97.2}{150.0}=0.648\text{ (64.8\%).}
\end{align}
Distance explains 64.8\% of delivery-time variability; the remaining 35.2\% is due to factors not included in the model (traffic, time of day, courier availability, and so on).
\end{solproblema}
